\noindent O esquema descreve cada signo usado nas tarefas pelo conteúdo e pela expressão. Não é o modelo de dados. A coluna \emph{Na observação do agente} aplica a equivalência de canal: \emph{trecho} é o texto geométrica e tipograficamente legível no recorte; abaixo do limiar de legibilidade restam o tipo do elemento e sua posição. Nenhum identificador interno, estado lógico de desabilitação ou elemento fora do recorte entra na observação. A coluna \emph{Evento} liga o signo ao registro do pesquisador. Os signos de aplicação recebem explicação completa nas três camadas da ajuda.
\newcommand{\SGhead}{\hline \textbf{Signo (entidade.atributo)} & \textbf{Origem} & \textbf{Conteúdo (tipo, restrições, default)} & \textbf{Expressão (widget)} & \textbf{Tipo} & \textbf{Na observação do agente} & \textbf{Evento}\\ \hline}
\begin{table*}[!t]
\caption{Esquema conceitual de signos estendido (parte 1: cardápio, item e navegação)}
\label{tab:signos1}
\centering\scriptsize\setlength{\tabcolsep}{3pt}\renewcommand{\arraystretch}{1.08}
\begin{tabularx}{\textwidth}{@{}>{\raggedright\arraybackslash}p{.135\textwidth}>{\raggedright\arraybackslash}p{.085\textwidth}>{\raggedright\arraybackslash}X>{\raggedright\arraybackslash}p{.15\textwidth}>{\raggedright\arraybackslash}p{.085\textwidth}>{\raggedright\arraybackslash}p{.165\textwidth}>{\raggedright\arraybackslash}p{.13\textwidth}@{}}
\SGhead
Restaurante.nome & domínio & texto; denominação do restaurante ou da lanchonete de treino & cabeçalho & estático & trecho legível & \UXev{elemento_visto}\\
Tela.assunto & aplicação & enumerado S0--S5; cena corrente & título da tela & estático & trecho do título & \UXev{tela_alterada}\\
Início.histórico & domínio & lista das linhas do último pedido, quando prevista na tarefa & relação de itens & estático & trechos de linhas & \UXev{elemento_visto}\\
Pedido.repetição & aplicação & ação; acrescenta as linhas históricas às atuais & botão Repetir pedido & estático & alvo visível & \UXev{item_adicionado}\\
Navegação.destino & convencional & enumerado: início, cardápio ou revisão & controles de navegação & estático & alvo visível & \UXev{tela_alterada}\\
Navegação.retorno & convencional & ação; abandona edição pendente & botão Voltar & estático & gesto ou alvo corrente & \UXev{tela_alterada}\\
Categoria.nome & domínio & texto; cinco categorias na ordem herdada & seção (A) ou aba (B) & estático & trecho & \UXev{elemento_visto}\\
Categoria.ativa & aplicação & enumerado; default: entradas (somente B) & realce da aba & dinâmico & pistas visíveis e itens expostos & \UXev{tela_alterada}\\
Item.nome & domínio & texto do catálogo; 30 itens & cabeçalho do cartão & estático & trecho, sem completar vizinhos & \UXev{elemento_visto}\\
Item.descrição & domínio & texto; ingredientes ou preparo & texto do cartão & estático & trecho & \UXev{elemento_visto}\\
Item.etiquetas & transformado & conjunto de \{V, SL, SA\}, possivelmente vazio & marcas junto ao item & metalinguístico & somente o próprio trecho & \UXev{elemento_visto}\\
Item.preço & domínio & centavos inteiros positivos; dois preços para massas & valor em reais & estático & trecho monetário & \UXev{elemento_visto}\\
Item.detalhes & aplicação & ação; abre a composição do item tocado & cartão tocável & estático & alvo visível, sem identificador interno & \UXev{tela_alterada}\\
Item.opção & domínio & tamanho ou acompanhamento obrigatório; default: nenhuma escolha & grupo de opções & dinâmico & textos expostos & \UXev{opcao_selecionada}\\
Opção.exigência & aplicação & condição para incluir o item & aviso escrito & metalinguístico & trecho do aviso, sem marca lógica de desabilitação & \UXev{elemento_visto}\\
Item.quantidade & domínio & inteiro $\geq 1$ & número e controles $+/-$ & dinâmico & número ou sinal visível & \UXev{acao_executada}\\
Item.inclusão & aplicação & ação; efetiva uma nova linha & botão Adicionar & dinâmico & alvo visível admissível & \UXev{item_adicionado}\\ \hline
\end{tabularx}
\end{table*}
\begin{table*}[!t]
\caption{Esquema conceitual de signos estendido (parte 2: pedido, ajuda, vista e protocolo)}
\label{tab:signos2}
\centering\scriptsize\setlength{\tabcolsep}{3pt}\renewcommand{\arraystretch}{1.08}
\begin{tabularx}{\textwidth}{@{}>{\raggedright\arraybackslash}p{.135\textwidth}>{\raggedright\arraybackslash}p{.085\textwidth}>{\raggedright\arraybackslash}X>{\raggedright\arraybackslash}p{.15\textwidth}>{\raggedright\arraybackslash}p{.085\textwidth}>{\raggedright\arraybackslash}p{.165\textwidth}>{\raggedright\arraybackslash}p{.13\textwidth}@{}}
\SGhead
Pedido.resumo & transformado & composição e total; somente em B & rodapé fixo & dinâmico & apenas a interseção efetiva com o recorte & \UXev{elemento_visto}\\
Resumo.página & aplicação & até quatro linhas por página; Anterior/Próxima quando há página nessa direção & paginação do rodapé (B) & dinâmico & alvos e linhas expostos & \UXev{acao_executada}\\
Linha.composição & transformado & nome, opção, unidades e subtotal & linha da revisão & dinâmico & fragmentos expostos & \UXev{elemento_visto}\\
Linha.unidades & transformado & inteiro positivo da linha existente & controles $+/-$ & dinâmico & controle visível & \UXev{item_alterado}\\
Linha.edição & aplicação & ação; reabre os detalhes da linha e aplica a mudança & botões Alterar e Salvar alteração & dinâmico & alvo do recorte & \UXev{item_alterado}\\
Linha.remoção & aplicação & ação; exclui a linha escolhida & botão Remover & dinâmico & alvo visível & \UXev{item_removido}\\
Linha.reposição & aplicação & ação; recupera a última remoção disponível & botão Desfazer remoção & dinâmico & alvo visível & \UXev{item_adicionado}\\
Pedido.total & transformado & soma dos subtotais, em reais & valor total & dinâmico & trecho & \UXev{elemento_visto}\\
Pedido.vazio & aplicação & ausência de linhas & mensagem & dinâmico & trecho & \UXev{elemento_visto}\\
Pedido.registro & aplicação & ação; exige comanda não vazia & botão Confirmar pedido, depois mensagem e linhas & dinâmico & alvo e resposta no recorte & \UXev{pedido_confirmado}\\
Ajuda.acesso & aplicação & ação; solicita explicação da cena & botão Ajuda ou F1 & estático & alvo visível ou acesso ubíquo previsto & \UXev{ajuda_aberta}\\
Ajuda.camada & aplicação & enumerado: dica local, sobre esta tela, guia de uso & três escolhas & dinâmico & opções da ajuda aberta & \UXev{ajuda_aberta}\\
Ajuda.tópico & aplicação & texto do assunto da cena de origem na camada escolhida & título e texto & metalinguístico & trecho, somente após a abertura & \UXev{observacao}\\
Ajuda.retorno & convencional & ação; retoma a cena de origem & botão Fechar ajuda ou Esc & estático & alvo visível & \UXev{ajuda_fechada}\\
Vista.escala e posição & convencional & janela de $360/z$ por $800/z$ px CSS e posição & gestos de ampliar, reduzir e deslocar & dinâmico & recorte e gestos permitidos & \UXev{zoom_alterado}, \UXev{vista_deslocada}\\
Protocolo.instrução & aplicação & enunciado literal e estável da tarefa & texto do protocolo & metalinguístico & texto literal & \UXev{tarefa_iniciada}\\
Protocolo.saída & convencional & ação; interromper por iniciativa própria & gesto Desistir & dinâmico & gesto de protocolo & \UXev{desistencia}\\ \hline
\end{tabularx}
\end{table*}
